% ===========================================================================
% RELATO DO ESTADO ATUAL DO TCC
% Documento curto, em formato de artigo, na mesma classe e no mesmo estilo da
% monografia (modelo TCC_PCS_EPUSP_2023).
% Compilar: latexmk relato.tex   (ou: make relato)
% ===========================================================================

\documentclass[article,12pt,oneside,a4paper,english,brazil]{abntex2}

\newif\ifparcial
\parcialtrue

% ---------------------------------------------------------------------------
% Preâmbulo: pacotes e macros compartilhados por parcial.tex e final.tex
% Compilado com XeLaTeX (ver latexmkrc), por causa da fonte Arial.
% ---------------------------------------------------------------------------

% ---------------------------------------------------------------------------
% Fontes. O texto é Arial 12 (o tamanho vem da opção 12pt da classe). Numa
% máquina sem Arial instalada, usa-se a TeX Gyre Heros, clone livre da
% Helvetica, de métrica equivalente. A capa mantém a fonte serifada com que foi
% desenhada (TeX Gyre Termes, clone livre da Times).
% ---------------------------------------------------------------------------
\usepackage{amsmath}
\usepackage{fontspec}
\IfFontExistsTF{Arial}{%
  \setmainfont{Arial}
  \setsansfont{Arial}
}{%
  \setmainfont{texgyreheros}[Extension=.otf, UprightFont=*-regular,
    BoldFont=*-bold, ItalicFont=*-italic, BoldItalicFont=*-bolditalic]
  \setsansfont{texgyreheros}[Extension=.otf, UprightFont=*-regular,
    BoldFont=*-bold, ItalicFont=*-italic, BoldItalicFont=*-bolditalic]
}
\newfontfamily\fontecapa{texgyretermes}[Extension=.otf, UprightFont=*-regular,
  BoldFont=*-bold, ItalicFont=*-italic, BoldItalicFont=*-bolditalic]
% Letras e dígitos das poucas fórmulas do texto na mesma fonte do texto
\usepackage[italic]{mathastext}

\usepackage{indentfirst}
% Folga para parágrafos com fórmulas ou termos longos que não se quebram
\setlength{\emergencystretch}{2em}
\usepackage{microtype}
\usepackage{graphicx}
\usepackage{xcolor}
\usepackage{booktabs}
\usepackage{tabularx}
% Coluna de parágrafo alinhada à esquerda, para tabelas com texto (evita a
% hifenização forçada das colunas estreitas justificadas)
\newcolumntype{L}[1]{>{\raggedright\arraybackslash}p{#1}}
\usepackage{listings}
\usepackage{tikz}
\usetikzlibrary{positioning,arrows.meta,fit,backgrounds}

% Citações no padrão ABNT, autor-data
\usepackage[brazilian,hyperpageref]{backref}
\usepackage[alf]{abntex2cite}

\renewcommand{\backrefpagesname}{Citado na(s) página(s):~}
\renewcommand{\backref}{}
\renewcommand*{\backrefalt}[4]{%
  \ifcase #1 %
    Nenhuma citação no texto.%
  \or
    Citado na página #2.%
  \else
    Citado #1 vezes nas páginas #2.%
  \fi}

% ---------------------------------------------------------------------------
% Estética: fora da capa, tudo em preto. As duas cores abaixo são usadas
% apenas pela capa (config/capa.tex).
% ---------------------------------------------------------------------------
\definecolor{azulpoli}{HTML}{0B2A4A}
\definecolor{cinzacapa}{HTML}{5A6270}

\renewcommand{\ABNTEXchapterfont}{\bfseries}
\renewcommand{\ABNTEXsectionfont}{\bfseries}
\renewcommand{\ABNTEXsubsectionfont}{\bfseries}

\AtBeginDocument{%
  \hypersetup{
    pdftitle={\imprimirtitulo},
    pdfauthor={\autorumnome, \autordoisnome},
    pdfsubject={\imprimirtipotrabalho},
    hidelinks,
  }%
}

% Código C nas listagens: sem cor e sem negrito
\lstset{
  language=C,
  basicstyle=\ttfamily\small,
  keywordstyle=,
  commentstyle=\itshape,
  numbers=left,
  numberstyle=\tiny,
  frame=single,
  breaklines=true,
  showstringspaces=false,
}

% ---------------------------------------------------------------------------
% Marcação de trabalho pendente.
% Para a versão de entrega, troque \mostrarpendentestrue por \mostrarpendentesfalse:
% os marcadores somem do PDF sem precisar apagá-los do texto.
% ---------------------------------------------------------------------------
\newif\ifmostrarpendentes
\mostrarpendentesfalse
\newcommand{\pendente}[1]{%
  \ifmostrarpendentes{\color{red}[PENDENTE: #1]}\fi}

% Sem negrito no corpo do texto: os rótulos de listas de descrição (QP1, nomes
% das métricas) saem em estilo normal, seguidos de dois-pontos.
\renewcommand*{\descriptionlabel}[1]{\hspace{\labelsep}\normalfont #1:}

% Termos recorrentes
\newcommand{\isr}{\textit{ISR}}
\newcommand{\llm}{\textit{LLM}}

% ---------------------------------------------------------------------------
% Dados do trabalho, no formato do modelo do PCS: alimentam a capa e a folha
% de rosto geradas pela classe abntex2.
% ---------------------------------------------------------------------------

\titulo{Detecção de Defeitos de Concorrência em Software Embarcado Orientado a
Interrupções, com Triagem e Reparo Assistidos por Modelos de Linguagem}

\autor{Lucas de Oliveira Martim \\ Pedro Henrique Machado Almeida}

\local{São Paulo, SP}
\data{2026}

\orientador{Prof. Dr. Carlos Cugnasca}

\instituicao{%
  Universidade de São Paulo -- USP
  \par
  Escola Politécnica
  \par
  Departamento de Engenharia de Computação e Sistemas Digitais (PCS)}

\tipotrabalho{Monografia de Conclusão de curso (Graduação)}


\titulo{Relato do Estado Atual do Trabalho de Conclusão de Curso}
\data{28 de setembro de 2026}
\tipotrabalho{Relato de andamento}
\preambulo{Relato de andamento do Trabalho de Conclusão de Curso apresentado ao
Departamento de Engenharia de Computação e Sistemas Digitais da Escola
Politécnica da Universidade de São Paulo.}

\begin{document}

\selectlanguage{brazil}
\frenchspacing

\imprimirfolhaderosto

\section{Identificação e objetivo}

O trabalho constrói uma ferramenta que detecta defeitos de concorrência em
código C embarcado orientado a interrupções e encerra o ciclo com um modelo de
linguagem, indo da detecção à contextualização, à triagem, ao reparo e à
reverificação. O compromisso central é a ausência de falsos negativos em
relação a uma lista explícita de premissas, aceitando falsos positivos em
troca, com o esforço de revisão controlado pela etapa de triagem em vez de pelo
descarte de candidatos. A fundamentação, os requisitos e o projeto estão na
monografia parcial entregue junto com este relato, que se limita a informar em
que ponto o trabalho se encontra.

O título completo é \emph{Detecção de Defeitos de Concorrência em Software
Embarcado Orientado a Interrupções, com Triagem e Reparo Assistidos por Modelos
de Linguagem}, sob orientação do Prof. Dr. Carlos Cugnasca.

\section{Pesquisa bibliográfica e decisões dela derivadas}

A fase de revisão está concluída. Foram analisados em profundidade onze artigos,
divididos em duas literaturas que não se comunicam. A primeira trata da detecção
de defeitos de concorrência em programas com interrupções e reúne quatro
ferramentas publicadas entre 2020 e 2026. A segunda trata do uso de modelos de
linguagem para apoiar analisadores estáticos e reúne sete trabalhos, todos
dedicados a defeitos sequenciais. A interseção entre as duas está vazia, e é ela
que o trabalho ocupa.

A revisão não se limitou à leitura. Cada artigo foi confrontado com os artefatos
correspondentes, quando obtíveis, e apenas um dos quatro detectores está
disponível, o que obrigou a tratar quase todo número publicado como afirmação do
autor. Desse confronto vieram dois resultados que a leitura isolada não daria. O
gabarito do \textit{benchmark} adotado foi recontado a partir das anotações do
próprio código e resultou em 48 defeitos e 38 armadilhas, ao passo que os três
artigos que o empregam usam contagens diferentes entre si. O conjunto de
programas reais distribuído com uma das ferramentas foi igualmente recontado e
mostrou-se corrompido em um de seus dezoito programas, o que explica a
divergência entre o número publicado pelos autores e o número registrado nos
arquivos. Ambas as verificações estão documentadas e são, elas próprias,
contribuição do trabalho, pois nenhuma taxa de acerto dessa literatura é
comparável enquanto o gabarito não estiver acordado.

Do trabalho de digestão saíram as decisões de projeto que hoje governam a
implementação. O recall tem precedência sobre a precisão, o que inverte a
prioridade adotada pela literatura. Nenhum candidato é descartado sem prova de
impossibilidade, de modo que o modelo de linguagem ordena e explica, mas nunca
remove. A unidade de contagem e a regra que decide quando um candidato
corresponde a um defeito anotado foram fixadas por escrito antes de qualquer
medição. E a afirmação de ausência de falsos negativos passou a ser acompanhada
da lista de premissas sob as quais vale, item que nenhum dos trabalhos revisados
apresenta.

\section{Planejamento}

O planejamento foi consolidado num roteiro de dez semanas, de 31 de agosto a 6
de novembro de 2026, organizado em duas frentes que avançam em paralelo e cinco
marcos de acompanhamento. A primeira frente cuida da análise estática, que
encontra os candidatos a defeito e reúne as evidências de cada um. A segunda
cuida da integração com modelos de linguagem, responsável pela triagem, pelo
reparo e pela avaliação dessas etapas. Uma terceira frente, o painel de
operação, é construída no tempo que as duas primeiras liberam.

O roteiro prevê que a segunda frente nunca dependa da primeira. Até a
integração, prevista para a oitava semana, ela trabalha sobre vinte casos de
teste montados manualmente a partir do \textit{benchmark}, o que permitiu que as
duas frentes avançassem desde a primeira semana sem bloqueio mútuo.

\section{O que já foi realizado}

\subsection{Análise estática}

O ambiente de compilação está operacional e todos os 31 programas do
\textit{benchmark} são processados pela ferramenta. O gabarito foi certificado
contra contagem manual, acompanhado de uma errata de doze anotações em que cada
correção traz sua evidência, sem que nenhuma delas altere as contagens.

A primeira das três etapas da análise está implementada e produz 566 candidatos
sobre os 31 programas, com média de 18,3 por programa. O teste que condiciona
todo o projeto passa, com os 48 defeitos anotados presentes entre os candidatos.

\subsection{Integração com modelos de linguagem}

O ambiente de execução está construído e permite trocar o modelo sem alterar as
instruções enviadas a ele, reaproveitar respostas já obtidas e registrar, para
cada medição, qual modelo e qual configuração a produziram. Os vinte casos de
teste estão montados e cobrem a variedade de situações que o \textit{benchmark}
apresenta, entre elas seis duplas adversariais, nas quais um defeito real e uma
armadilha diferem em um único fato.

A sonda de viabilidade, que era o item de maior risco do projeto por não ter
precedente na literatura, foi respondida positivamente. Sem nenhuma das técnicas
previstas para a etapa de triagem, nenhum defeito foi classificado como inviável
e as duplas adversariais foram separadas corretamente. Duas das cinco linhas da
tabela de ablação já estão medidas, conforme o Quadro~\ref{quadro:ablacao}.

\begin{quadro}[htb]
  \caption{\label{quadro:ablacao}Linhas da ablação já medidas}
  \small
  \begin{tabular}{|l|L{2.6cm}|L{2.8cm}|L{3.0cm}|}
    \hline
    Linha & Recall & Rejeição de armadilhas & \textit{Inspection Ratio} \\
    \hline
    instrução simples & 100\% (10/10) & 70\% & 65\% \\
    \hline
    mais regras do domínio & 100\% (10/10) & 50\% & 75\% \\
    \hline
  \end{tabular}
  \fonte{Elaborado pelos autores. Medições sobre 20 casos de teste, uma amostra
  por candidato, setembro de 2026.}
\end{quadro}

O resultado da segunda linha é negativo e está sendo tratado como achado, e não
como ruído a corrigir. Ensinar ao modelo as regras do domínio piorou seu
julgamento sobre duas armadilhas, sem mover o recall. Com vinte casos e uma
única amostra por candidato, a evidência ainda é fraca, e é a repetição das
medições que permitirá afirmar a diferença. As demais etapas previstas para essa
frente, entre elas o reparo, estão implementadas e aguardam medição.

\subsection{Produtos disponíveis}

Estão disponíveis o código das duas frentes, acompanhado de mais de duzentos
testes automatizados, a base de conhecimento com a revisão da literatura e um
diário de implementação em português, com uma página por semana de trabalho, que
registra o que foi construído e por quê.

\section{Aderência ao cronograma}

Os dois primeiros marcos foram entregues, conforme o Quadro~\ref{quadro:marcos}.
O trabalho está na quinta das dez semanas previstas.

\begin{quadro}[htb]
  \caption{\label{quadro:marcos}Situação dos marcos}
  \small
  \begin{tabular}{|l|l|L{6.8cm}|l|}
    \hline
    Marco & Data & Entrega & Situação \\
    \hline
    M1 & 11/09 & gabarito certificado, casos de teste montados e sonda de viabilidade respondida & concluído \\
    \hline
    M2 & 25/09 & primeira etapa da análise nos 31 programas, com os 48 defeitos encontrados & concluído \\
    \hline
    M3 & 02/10 & tabela de ablação da triagem & em curso \\
    \hline
    M4 & 23/10 & ciclo completo no \textit{benchmark}, da análise ao reparo & previsto \\
    \hline
    M5 & 06/11 & suíte de programas reais e métricas finais & previsto \\
    \hline
  \end{tabular}
  \fonte{Elaborado pelos autores.}
\end{quadro}

O marco M2 foi entregue uma semana antes da data prevista. O marco M1 teve um de
seus três itens concluído com atraso, pois a sonda de viabilidade dependia de
acesso a um provedor de modelos de linguagem, obtido apenas em 26 de setembro.
Os outros dois itens saíram na data, e a frente afetada usou o período para
construir as partes que não dependiam desse acesso, de modo que o atraso não se
propagou para os marcos seguintes.

\section{Próximas etapas}

Na análise estática, as próximas três semanas implementam as duas etapas
restantes, que reduzem o conjunto de candidatos descartando apenas aqueles cuja
impossibilidade pode ser provada. A oitava semana integra as duas frentes, a
nona aplica a ferramenta à suíte de 18 programas reais e a décima consolida os
resultados e o texto final.

Na frente de modelos de linguagem, as três linhas restantes da ablação precisam
ser medidas num provedor com capacidade adequada, uma vez que o serviço gratuito
usado até aqui limita a vazão a ponto de tornar a medição completa inviável em
tempo. O custo foi levantado e é da ordem de dez dólares para a tabela completa.
Uma dessas linhas foi deliberadamente adiada para depois da integração, quando
passará a existir contexto real a ser consultado, sem o qual ela não mediria
aquilo que lhe dá nome. Em seguida vêm a medição das taxas de reparo e a
repetição da triagem sobre os candidatos produzidos pela análise, que formam
população diferente e mais difícil que a dos casos montados manualmente.

\end{document}
